\subsection{Understanding deletion, link, and the degree shift}
Before proving the block decomposition, consider a filled triangle on vertices $c,u,v$. Deleting $c$ leaves the edge $uv$ and its vertices. The link of $c$ also contains that edge, because adding $c$ to $uv$ gives a triangle in the original complex. If only the three boundary edges are present, the deletion still contains $uv$, but the link consists of the two isolated vertices $u,v$: adding $c$ to the edge no longer gives an included cell. Thus the deletion asks what remains without the center, whereas the link asks which cells participate in cells incident to the center.

The word ``local'' here is combinatorial. The entire complex has already been restricted to a geometric neighborhood. The link further selects the pattern of cells attached to its distinguished vertex. It should not be confused with another distance-cutoff neighborhood, nor with the complete subgraph on the center's neighbors. A link edge requires the corresponding center-containing triangle. The two center edges alone are insufficient.

Every edge containing $c$ corresponds to a vertex of its link. Every triangle containing $c$ corresponds to an edge of the link. This accounts for the shift by one degree: center-containing degree-one cochains look like degree-zero link cochains. The center vertex itself needs a counterpart one degree below link vertices. That extra coordinate is the \emph{augmentation}, formally associated with an empty simplex of degree $-1$.

For ordinary augmented cochains, a scalar in degree $-1$ maps to the same constant on every link vertex. Dividing degree-zero cohomology by these constants leaves differences between connected components. If the link has $k>0$ components, this reduced degree-zero dimension is $k-1$. In the geometric sheaf, the augmentation map is weighted, but an invertible cochain rescaling gives the same reduced dimension. The weights still affect positive eigenvalues.

The triangle outline provides a useful check. Its deletion is a single edge, with no one-dimensional hole. Its link has two components, so it has one independent component difference. The center-zero degree-one nullity is therefore one. Filling the triangle makes the link connected and removes that contribution. In more complicated neighborhoods, deletion cycles and reduced link components can both contribute. A degree-one zero mode in this sheaf need not describe only a hole of the original unweighted complex.
